%% ── 8. DESARROLLO Y CONSTRUCCIÓN ─────────────────────────────────────────────
\chapter{Desarrollo y Construcción}

Este capítulo documenta el proceso de desarrollo de ContractIA, desde los
primeros experimentos en notebooks hasta el MVP desplegado en producción.
El desarrollo siguió un enfoque iterativo e incremental, con 15 versiones
de notebook (prototipo) y 4 versiones de plataforma (MVP), priorizando en
cada ciclo la hipótesis de mayor riesgo técnico pendiente de validar.

\section{Enfoque Metodológico}

El desarrollo adoptó un ciclo corto de construcción–medición–aprendizaje
(\textit{build–measure–learn}), con las siguientes decisiones de proceso:

\begin{itemize}
    \item \textbf{Iteraciones cortas (2--5 días):} cada versión del notebook
    resolvía un problema específico antes de agregar nueva funcionalidad.
    \item \textbf{Contrato piloto fijo:} el contrato A
    (324 páginas) fue el benchmark constante entre iteraciones, permitiendo
    comparaciones objetivas de cobertura, tiempo y hallazgos.
    \item \textbf{Separación prototipo / MVP:} las iteraciones v1--v15 se
    desarrollaron en Jupyter como exploración analítica; la plataforma
    productiva (FastAPI + Cloud Run) se construyó en paralelo a partir de v9.5.
    \item \textbf{Validación continua:} tras cada iteración se revisó
    manualmente una muestra de 10--20 hallazgos para detectar regresiones.
\end{itemize}

\section{Historial de Versiones}

\begin{table}[H]
\centering
\caption{Evolución de versiones del sistema ContractIA}
\label{tab:versiones}
\begin{tabular}{p{2cm} p{3.5cm} p{7cm}}
\toprule
\textbf{Versión} & \textbf{Hito} & \textbf{Cambio principal} \\
\midrule
v1.0  & Pipeline base      & Segmentación PDF + FAISS + 1 agente LLM \\
v2--v4 & Calidad RAG       & Chunking adaptativo; filtrado anti-TOC \\
v5    & Multi-agente       & Jurista + Auditor + Cronista con Pydantic \\
v6--v7 & Orquestación      & \texttt{asyncio.gather}; trazabilidad por cláusula \\
v8.7  & Hybrid RAG         & BM25 + FAISS fusión RRF \\
v8.8  & Cohere Reranker    & Cross-encoder sobre top-20 candidatos \\
v9.0  & GraphRAG           & DiGraph por sección; extracción de tripletas LLM \\
v9.5  & API + Cloud Run    & FastAPI; Docker; CI/CD con GitHub Actions \\
v9.6  & Telegram bot       & Webhook; flujo de aprobación de usuarios \\
v9.7  & Frontend Next.js   & Panel web; polling de estado en tiempo real \\
v9.8  & Cola + notif.      & Cola por usuario (máx. 5 trabajos); email al finalizar \\
\bottomrule
\end{tabular}
\end{table}

\section{Iteraciones del Desarrollo}

\subsection{Iteración 1 — Prototipo en Notebook (v1.0--v4)}

\textbf{Objetivo:} validar que un LLM puede identificar inconsistencias
contractuales en un contrato de concesión peruano real.

\textbf{Construcción:}
\begin{itemize}
    \item Extracción de texto con PyMuPDF y segmentación por jerarquía
    de capítulos y cláusulas (regex + heurísticas).
    \item Indexación de fragmentos en FAISS con embeddings
    \texttt{text-embedding-004} de Vertex AI.
    \item Un único agente LLM (Gemini 2.5 Pro) con prompt de análisis
    semántico e instrucción de devolver JSON estructurado.
    \item Iteraciones v2--v4: mejora del chunking adaptativo para evitar
    que el índice de contenido (TOC) sea indexado como cláusula.
\end{itemize}

\textbf{Desafío clave:} el módulo anti-TOC era crítico; sin él, el RAG
recuperaba fragmentos del índice del contrato en lugar de cláusulas reales,
generando hallazgos sin sustento.

\textbf{Resultado:} 35 inconsistencias detectadas sobre el Contrato A
(314 páginas) en 47 minutos. Validó la hipótesis central de viabilidad técnica.

\subsection{Iteración 2 — Arquitectura Multi-Agente (v5--v7)}

\textbf{Objetivo:} aumentar la cobertura y especialización del análisis
descomponiendo el agente único en roles diferenciados.

\textbf{Construcción:}
\begin{itemize}
    \item \textbf{Jurista:} analiza consistencia con normativa (D.Leg.\ 1362,
    guías ProInversión). Devuelve \texttt{HallazgoJurista} con campo
    \texttt{referencia\_normativa}.
    \item \textbf{Auditor:} detecta contradicciones internas, referencias
    cruzadas rotas y omisiones entre secciones. Devuelve
    \texttt{HallazgoAuditor} con \texttt{seccion\_referenciada}.
    \item \textbf{Cronista:} identifica conflictos de fechas, plazos
    inconsistentes y hitos mal ordenados. Devuelve \texttt{HallazgoCronista}
    con \texttt{fecha\_detectada} y \texttt{fecha\_esperada}.
    \item Orquestador con \texttt{asyncio.gather} para ejecutar los tres
    agentes en paralelo por sección, reduciendo la latencia total.
    \item Esquemas de salida con Pydantic v2 para garantizar estructura
    consistente y validación automática de tipos.
\end{itemize}

\textbf{Desafío clave:} alinear el formato de salida JSON de tres agentes
distintos en un único consolidado por sección, evitando duplicados y
asignando correctamente el índice global de hallazgo.

\textbf{Resultado:} incremento de cobertura temática; hallazgos clasificados
por tipo (jurídico, estructural, temporal) con trazabilidad completa.

\subsection{Iteración 3 — RAG Híbrido y GraphRAG (v8--v9)}

\textbf{Objetivo:} mejorar la relevancia del contexto normativo recuperado
y enriquecer el análisis con relaciones entre secciones del contrato.

\textbf{Construcción:}
\begin{itemize}
    \item \textbf{Hybrid RAG (v8.7):} fusión de BM25 (recuperación léxica)
    y FAISS (recuperación semántica) mediante Reciprocal Rank Fusion (RRF),
    mejorando la cobertura sobre términos legales técnicos poco frecuentes.
    \item \textbf{Cohere Reranker (v8.8):} cross-encoder sobre los
    top-20 candidatos del RAG híbrido para seleccionar los 5 más relevantes,
    reduciendo ruido en el contexto entregado a los agentes.
    \item \textbf{GraphRAG (v9.0):} extracción de tripletas
    \textit{(sujeto, relación, objeto)} por sección mediante LLM;
    construcción de un \texttt{DiGraph} (NetworkX) por contrato con
    relaciones: \texttt{obliga\_a}, \texttt{referencia}, \texttt{modifica},
    \texttt{contradice}, \texttt{complementa}. Permite consultar vecinos
    de un nodo para enriquecer el contexto de análisis cruzado.
\end{itemize}

\textbf{Desafío clave:} la extracción de tripletas con LLM es costosa
(1 llamada adicional por sección); se implementó caché por hash de sección
para evitar reextracción en análisis sucesivos del mismo contrato.

\textbf{Resultado:} reducción de falsos positivos por contexto irrelevante;
detección de contradicciones entre secciones no adyacentes gracias al grafo.

\subsection{Iteración 4 — Despliegue en Producción (v9.5--v9.8)}

\textbf{Objetivo:} convertir el notebook analítico en un servicio accesible,
operativo y escalable para usuarios no técnicos.

\textbf{Construcción:}
\begin{itemize}
    \item \textbf{API FastAPI (v9.5):} encapsulamiento del pipeline de
    auditoría como servicio HTTP con endpoints \texttt{/contracts/upload},
    \texttt{/contracts/\{id\}/audit} y \texttt{/contracts/\{id\}/status}.
    Containerización con Docker; imagen publicada en Artifact Registry.
    CI/CD con GitHub Actions: build → push → deploy a Cloud Run en cada
    merge a \texttt{main}.
    \item \textbf{Bot Telegram (v9.6):} webhook en \texttt{/telegram/webhook};
    flujo de registro, aprobación por administrador (rol \texttt{auditor})
    y comandos \texttt{/upload}, \texttt{/status}, \texttt{/report}.
    \item \textbf{Frontend Next.js (v9.7):} panel web en
    \texttt{contractia.pe}; autenticación JWT; subida de contratos;
    visualización de hallazgos con filtro por tipo y severidad;
    polling de estado en tiempo real (actualización cada 10 segundos).
    \item \textbf{Cola y notificación (v9.8):} sistema de cola por usuario
    con máximo 5 trabajos pendientes; procesamiento secuencial en segundo
    plano; notificación por correo electrónico al finalizar el análisis.
\end{itemize}

\textbf{Desafío clave:} garantizar que el análisis de un contrato extenso
(análisis completo: $\approx 75$ minutos) no bloqueara el hilo principal
de la API; solución: \texttt{BackgroundTasks} de FastAPI con persistencia
de estado en Cloud SQL.

\textbf{Resultado:} MVP funcional en producción en \texttt{contractia.pe},
accesible vía web y Telegram, con análisis completo del Contrato A
en $\approx 1$ hora 15 minutos con Gemini~2.5~Pro.

\section{Dataset Utilizado}

\begin{table}[H]
\centering
\caption{Datasets empleados en el desarrollo y validación del MVP}
\label{tab:datasets}
\begin{tabular}{p{3.5cm} p{9.5cm}}
\toprule
\textbf{Dataset} & \textbf{Descripción} \\
\midrule
Contrato piloto &
    Contrato A (ProInversión): 324 páginas, 41 secciones,
    448 cláusulas, 749{,}499 caracteres. Usado como benchmark
    constante entre iteraciones. \\
Corpus de segmentación &
    17 contratos de concesión de ProInversión de distintos sectores
    (agua, energía, vial), extensión variable. Usados exclusivamente
    para pruebas de robustez del módulo anti-TOC y segmentación. \\
Base normativa &
    Guías institucionales de ProInversión, D.Leg.\ 1362 y su
    Reglamento, normativa sectorial complementaria; vectorizados
    en FAISS e indexados en BM25. Total: $\approx 1{,}200$ fragmentos
    normativos. \\
\bottomrule
\end{tabular}
\end{table}

\section{Modelos Implementados y Evaluados}

\begin{table}[H]
\centering
\caption{Modelos LLM evaluados en ContractIA}
\label{tab:modelos}
\begin{tabular}{p{3.8cm} p{2.5cm} p{6.5cm}}
\toprule
\textbf{Modelo} & \textbf{Estado} & \textbf{Observación} \\
\midrule
Gemini 2.5 Pro (Vertex AI) & \textbf{Producción} &
    Velocidad óptima ($\approx 75$ min para 41 secciones); 119
    hallazgos sobre Contrato A; balance costo/cobertura. \\
Gemini 3.1 Pro Preview      & Evaluación &
    Mayor riqueza del grafo de conocimiento; 55 hallazgos en
    benchmark preliminar; mayor latencia ($\approx 3$ h 10 min). \\
Claude Sonnet 4.6 / Opus 4.6 & Pendiente &
    Integrado vía Vertex AI; en evaluación de cuota; resultados
    preliminares prometedores en análisis semántico. \\
\bottomrule
\end{tabular}
\end{table}

\section{Desafíos Técnicos y Soluciones}

Durante el desarrollo se enfrentaron tres desafíos técnicos principales:

\begin{enumerate}
    \item \textbf{Variabilidad en la estructura del PDF:}
    Los contratos de ProInversión usan numeraciones inconsistentes,
    tablas anidadas y TOC de múltiples páginas. Solución: módulo anti-TOC
    basado en detección estadística de densidad de referencias numéricas;
    heurística de longitud mínima de fragmento ($> 200$ caracteres) para
    filtrar entradas del índice.

    \item \textbf{Límite de contexto de los LLMs:}
    Cláusulas muy extensas exceden la ventana de contexto óptima.
    Solución: truncamiento adaptativo al 80\% del límite del modelo,
    preservando la cabecera de la cláusula y el párrafo más relevante
    según el score del Reranker.

    \item \textbf{Consistencia entre ejecuciones:}
    Los LLMs generativos producen variaciones entre ejecuciones con el
    mismo input. Solución: temperatura fija a 0.1 en todos los agentes;
    esquemas Pydantic con validación estricta de tipos; postprocesamiento
    de deduplicación por similitud de texto ($> 0.9$ cosine similarity).

    \item \textbf{Gestión de secretos y seguridad en producción:}
    Credenciales de API (Vertex AI, Cohere, PostgreSQL, SMTP) no pueden
    residir en el código ni en el Dockerfile. Solución: todas las
    credenciales almacenadas en Google Cloud Secret Manager; el servicio
    las lee en tiempo de ejecución mediante la service account
    \texttt{contractia-sa}; el Dockerfile no contiene ningún secreto.

    \item \textbf{Análisis de larga duración sin bloquear la API:}
    Un análisis completo toma $\approx 75$ minutos; un endpoint síncrono
    agotaría el timeout de Cloud Run (máx. 60 min por defecto) y bloquearía
    la respuesta al usuario. Solución: endpoint \texttt{/audit} dispara un
    \texttt{BackgroundTask} y devuelve inmediatamente el \texttt{job\_id};
    el cliente consulta el estado vía polling en \texttt{/status}; al
    finalizar el backend envía notificación por correo electrónico.
\end{enumerate}

\section{Estructura del Repositorio}

El código del MVP está organizado en dos capas principales: la lógica analítica
(\texttt{contractia/}) y la capa de servicio (\texttt{api/}), con separación
clara de responsabilidades.

\begin{longtable}{p{5cm} p{9cm}}
\caption{Módulos principales del repositorio ContractIA}
\label{tab:repo} \\
\toprule
\textbf{Módulo / Archivo} & \textbf{Responsabilidad} \\
\midrule
\endfirsthead
\multicolumn{2}{l}{\small\itshape Continuación de la Tabla~\ref{tab:repo}} \\[4pt]
\toprule
\textbf{Módulo / Archivo} & \textbf{Responsabilidad} \\
\midrule
\endhead
\midrule
\multicolumn{2}{r}{\small\itshape Continúa en la siguiente página} \\
\endfoot
\bottomrule
\endlastfoot
\texttt{contractia/config.py} &
    Variables de entorno, flags de feature (\texttt{RAG\_ENABLED},
    \texttt{GRAPHRAG\_ENABLED}, \texttt{AGENTIC\_RAG\_ENABLED}),
    modelos LLM configurables. \\[6pt]
\texttt{contractia/pipeline.py} &
    Orquestación de las fases FASE~0 a FASE~5: segmentación,
    indexación, RAG, agentes, consolidación, reporte. \\[6pt]
\texttt{contractia/orchestrator.py} &
    Función \texttt{auditar\_consistencia()} por sección;
    gestión de contexto RAG y llamadas a agentes. \\[6pt]
\texttt{contractia/agents/} &
    Fábrica de agentes (\texttt{factory.py}); clases
    \texttt{AgenteJurista}, \texttt{AgenteAuditor},
    \texttt{AgenteCronista}; schemas Pydantic por agente. \\[6pt]
\texttt{contractia/rag/} &
    Módulos FAISS, BM25, Cohere Reranker, GraphRAG (NetworkX);
    función unificada \texttt{recuperar\_contexto()}. \\[6pt]
\texttt{contractia/segmentation/} &
    Extracción PDF (PyMuPDF + pdfminer), anti-TOC,
    segmentador jerárquico, normalización de texto. \\[6pt]
\texttt{contractia/reports/} &
    Generación de reporte ejecutivo PDF (FPDF2);
    consolidación y deduplicación de hallazgos. \\[6pt]
\texttt{api/main.py} &
    Entrada FastAPI; registro de routers;
    endpoint de salud \texttt{/health}. \\[6pt]
\texttt{api/routers/} &
    \texttt{auth\_router} (JWT), \texttt{contracts\_router}
    (upload, audit, status, report), \texttt{admin\_router}
    (gestión de usuarios y roles). \\[6pt]
\texttt{contractia/telegram/} &
    Handler del bot; comandos; flujo de aprobación;
    wrapper psycopg2 para Cloud SQL. \\[6pt]
\texttt{Dockerfile} &
    Imagen Python 3.11 slim; sin secretos embebidos;
    \texttt{ENTRYPOINT} Uvicorn. \\[6pt]
\texttt{.github/workflows/deploy.yml} &
    Pipeline CI/CD: build → push a Artifact Registry →
    deploy a Cloud Run (Workload Identity Federation). \\
\end{longtable}

\section{Infraestructura y Pipeline CI/CD}

La infraestructura del MVP está desplegada íntegramente en Google Cloud
Platform, organizada en los siguientes servicios:

\begin{itemize}
    \item \textbf{Cloud Run:} servicio sin servidor (\textit{serverless})
    que ejecuta el contenedor Docker del API; escala automáticamente a
    cero cuando no hay tráfico, reduciendo costos operativos.
    \item \textbf{Cloud SQL (PostgreSQL 15):} base de datos relacional
    para usuarios, roles, jobs de análisis y metadatos de contratos.
    Acceso desde Cloud Run mediante el conector oficial de Cloud SQL.
    \item \textbf{Cloud Storage:} almacenamiento de contratos PDF subidos
    por los usuarios y reportes generados; acceso con signed URLs de
    duración limitada para descarga segura.
    \item \textbf{Artifact Registry:} registro de imágenes Docker;
    versiones inmutables por cada build del CI/CD.
    \item \textbf{Secret Manager:} almacén centralizado de credenciales
    (Vertex AI, Cohere, PostgreSQL, SMTP, JWT secret); sin secretos
    en variables de entorno del Dockerfile.
    \item \textbf{Vertex AI:} API administrada de Google para Gemini 2.5
    Pro (inferencia LLM) y \texttt{text-embedding-004} (embeddings);
    facturación por token, sin infraestructura de GPU propia.
\end{itemize}

El pipeline CI/CD ejecuta los siguientes pasos en cada \textit{merge} a la
rama \texttt{main}:

\begin{enumerate}
    \item \textbf{Checkout y autenticación:} obtención del token de acceso
    a GCP mediante Workload Identity Federation (sin clave de servicio
    almacenada en GitHub Secrets).
    \item \textbf{Build Docker:} construcción de la imagen con tag
    \texttt{git SHA} para trazabilidad de versiones.
    \item \textbf{Push a Artifact Registry:} imagen disponible en
    \texttt{us-central1-docker.pkg.dev/striking-domain-488917-p8/contractia/api}.
    \item \textbf{Deploy a Cloud Run:} actualización del servicio
    \texttt{contractia-api} con la nueva imagen; Cloud Run realiza
    \textit{rolling deployment} sin downtime.
    \item \textbf{Verificación de salud:} el workflow consulta
    \texttt{/health} y falla el deploy si no responde HTTP~200.
\end{enumerate}

\section{Estrategia de Pruebas y Validación}

La validación del sistema se realizó en tres niveles complementarios:

\begin{enumerate}
    \item \textbf{Validación funcional por iteración:}
    Tras cada versión de notebook se ejecutó el análisis completo sobre
    el contrato A y se comparó el conteo de hallazgos,
    tipos detectados y tiempo de procesamiento con la versión anterior.
    Cualquier regresión ($> 10\%$ de reducción en hallazgos) bloqueaba
    el avance a la siguiente iteración.

    \item \textbf{Validación manual de calidad semántica:}
    Una muestra aleatoria de 30 hallazgos de severidad ALTA fue revisada
    por dos especialistas legales independientes. Cada hallazgo fue
    clasificado como verdadero positivo (TP), falso positivo (FP) o
    falso negativo (FN). La precisión estimada resultante fue $\approx 87\%$.

    \item \textbf{Pruebas de integración del MVP:}
    Antes del despliegue en producción se verificaron manualmente:
    \begin{itemize}
        \item Subida de PDF vía API y vía bot Telegram.
        \item Flujo completo de análisis en segundo plano y consulta de estado.
        \item Generación y descarga del reporte ejecutivo PDF.
        \item Flujo de aprobación de usuario y cambio de rol.
        \item Cola de trabajos: llegada al límite de 5 y rechazo del 6.º.
        \item Notificación por correo al finalizar el análisis.
    \end{itemize}
\end{enumerate}

\begin{table}[H]
\centering
\caption{Resumen de resultados de validación por iteración}
\label{tab:validacion_iter}
\begin{tabular}{p{3cm} p{2.8cm} p{2.2cm} p{4.5cm}}
\toprule
\textbf{Versión} & \textbf{Hallazgos} & \textbf{Tiempo} & \textbf{Mejora clave} \\
\midrule
v1.0 (1 agente)      & 35  & 47 min  & Baseline del prototipo \\
v5   (multi-agente)  & 62  & 58 min  & $+77\%$ hallazgos \\
v8.8 (Hybrid RAG)    & 89  & 65 min  & Reducción FP; mejor contexto \\
v9.0 (GraphRAG)      & 103 & 72 min  & Detección cruzada entre secciones \\
v9.8 (MVP producción) & 119 & $\approx 75$ min & Benchmark definitivo en producción \\
\bottomrule
\end{tabular}
\end{table}

\section{Decisiones de Implementación}

\begin{table}[H]
\centering
\caption{Decisiones tecnológicas clave del MVP y su justificación}
\label{tab:decisiones_impl}
\begin{tabular}{p{3.5cm} p{4.5cm} p{5cm}}
\toprule
\textbf{Decisión} & \textbf{Alternativa considerada} & \textbf{Justificación} \\
\midrule
FastAPI sobre Flask &
    Flask, Django REST &
    Soporte nativo async/await; validación automática con Pydantic;
    OpenAPI generado automáticamente. \\
Cloud Run sobre GKE &
    GKE, Cloud Functions &
    Sin gestión de clúster; escala a cero; costo proporcional al uso;
    adecuado para carga irregular. \\
PostgreSQL sobre Firestore &
    Firestore, BigQuery &
    Esquema relacional para usuarios/roles/jobs; transacciones ACID;
    experiencia del equipo. \\
BackgroundTasks sobre Celery &
    Celery + Redis, Cloud Tasks &
    Sin infraestructura adicional para el MVP; suficiente para
    un análisis concurrente por usuario. \\
Polling sobre WebSockets &
    SSE, WebSockets &
    Compatibilidad universal con proxies y Telegram;
    implementación sin estado en Cloud Run. \\
Vertex AI sobre OpenAI API &
    OpenAI, Azure OpenAI &
    Residencia de datos en GCP; integración nativa con el resto de
    la infraestructura; facturación unificada. \\
\bottomrule
\end{tabular}
\end{table}
